%%%%%%%%%%  scan idx 23  |  folio 10  %%%%%%%%%%
de
Modules inversibles amples sur $X$ et $Y$.

\subsubsection*{4.1.}
Notons d'abord que \uline{pour un sch\'ema $X$ (ou plus g\'en\'eralement, un topos localement annel\'e) quelconque, il semble hors de doute que le ``bon'' invariant qui doit remplacer l'anneau $K(X)$ de [RRR]} n'est pas \uline{l'anneau $K^\bullet(X)$ consid\'er\'e pr\'ec\'edemment, d\'efini en termes de Modules localement libres sur $X$, mais bien le groupe $K(\Parf(X))$, envisag\'e dans 3.1., d\'efini en termes de complexes parfaits sur $X$, groupe que nous noterons ici} $K^\bullet(X)$, et que nous distinguerons du groupe ``naif'' de 3.1, que nous notons $K^\bullet_{\mathrm{naif}}(X)$. Ce point de vue est impos\'e par la n\'ecessit\'e de d\'efinir un homomorphisme (3.1) image directe pour un morphisme $f:X\to Y$ propre et de tor-dimension finie, de sch\'emas noeth\'eriens disons (le cas non noeth\'erien exigeant des propri\'et\'es de finitude locales un peu plus fortes sur $f$, cf.\ Exp.\ III). Avec la d\'efinition de $K^\bullet(X)$ adopt\'ee maintenant, cet homomorphisme $f_*$ se d\'efinit en effet de fa\c{c}on triviale par la formule (3.2), o\`u maintenant $\underline F$ est un complexe parfait quelconque sur $X$, ce qui implique que $Rf_*(\underline F)$ est \'egalement parfait.

\subsubsection*{4.2}
Cette nouvelle d\'efinition de $K^\bullet(X)$ soul\`eve cependant un nouveau probl\`eme d'alg\'ebre homologique, qui est celui d'une d\'efinition d'une $\lambda$-structure sur $K^\bullet(X)$, compatible avec sa structure d'anneau naturelle provenant du produit tensoriel dans la cat\'egorie d\'eriv\'ee. (NB Si $L_\bullet$ et $L'_\bullet$ sont deux complexes parfaits, il en est de m\^eme de leur produit tensoriel total $L_\bullet\otimes^{\mathbf L}L'_\bullet$). Plus pr\'ecisement, il y aurait lieu de d\'efinir des op\'erations ``puissances ext\'erieures'' dans la cat\'egorie d\'eriv\'ee $D^-(X)$, transformant complexes parfaits en complexes parfaits. Cette question n'a pas encore \'et\'e tir\'ee au clair \`a l'heure actuelle, et comme il a \'et\'e dit dans l'Introduction, elle constitue la lacune la plus s\'erieuse dans le S\'eminaire. Un fois cette question r\'esolue, et proc\'edant comme dans 2.3, on trouve une filtration naturelle sur $K^\bullet(X)$, permettant de former l'anneau gradu\'e associ\'e $\Gr^\bullet(X)$, qui jouera le r\^ole de l'anneau de Chow, et de d\'efinir une th\'eorie des classes de Chern
\[
c_i:K^\bullet(X)\longrightarrow \Gr^\bullet(X).
\]

\subsubsection*{4.3}
Pour donner un sens \`a la formule (1.2) pour un morphisme propre localement d'intersection compl\`ete $f:X\to Y$ de sch\'emas noeth\'eriens, il faut encore inclure dans l'\'enonc\'e de cette formule les inclusions non triviales (3.4) (o\`u $d$ est la dimension relative virtuelle) pour permettre de d\'efinir (3.3), et de d\'efinir enfin la classe tangente relative virtuelle
